\section{5. Conclusion}

\section{5. Conclusion}

In this paper, we propose a novel Unified Dual-Mask Physical Model (UDMPM) for robust light field depth estimation from spatially-varying non-Lambertian scenes. To address the core requirements for overcoming the failure of traditional epipolar plane image (EPI) assumptions under specular and glossy reflections, this work constructs a composite light field benchmark, proposes and validates UDMPM – a unified framework for physically grounded depth estimation. By fundamentally rethinking the photometric formation process, our approach departs from the restrictive single-Lambertian assumption, offering a theoretically sound solution to the long-standing "texture copying" degradation in non-Lambertian regions. A series of experiments conducted on diverse datasets validate the superiority of UDMPM across multiple dimensions, particularly in mitigating artifacts in mixed and challenging reflective scenarios. 

Specifically, the core contributions of this work are summarized as follows:
*   \textbf{Unified Dual-Mask Physical Modeling Mechanism:} We establish a unified theoretical framework capable of accommodating pure diffuse, specular/scattering, and urban mixed scenes. By modulating multi-directional epipolar geometries with dual masks, this mechanism fundamentally explains the physical root causes of depth estimation failure in non-Lambertian regions, breaking the limitations of traditional single-Lambertian assumptions.
*   \textbf{MRI-Inspired Pixel-Level Angular Frequency Material Parsing:} We ingeniously cross-apply frequency-domain analysis from medical imaging to computational photography. This method achieves real-time ($\mathcal{O}(HW)$ complexity) pixel-level material classification and BRDF parameter estimation without relying on heavy neural networks, effectively resolving the fine-grained material recognition bottleneck under $9 \times 9$ sparse angular sampling where purely data-driven CNNs fail.
*   \textbf{Component-Aware Adaptive Depth Estimation and Fusion Strategy:} We design an adaptive fusion mechanism that dynamically weights depth cues based on the parsed material components. This strategy effectively eliminates the "texture copying" failure caused by the violation of view-consistency assumptions in traditional EPI methods for high-reflective and transparent materials, significantly enhancing depth accuracy and robustness in both mixed and pure non-Lambertian scenarios.

The current validation is based on a fixed $9 \times 9$ sparse angular sampling and reveals certain inherent limitations. Specifically, complex textures in Lambertian regions still cause EPI slope blurring, and the mean absolute error (MAE) in extreme non-Lambertian scenarios remains constrained by the physical ambiguity of sparse views. Furthermore, the defocus branch struggles to learn effective depth features independently to complement the EPI cues. Finally, since we showed that merely increasing the number of views does not fundamentally resolve the non-Lambertian ambiguity, future work includes deriving an ambiguity-space resolution mechanism using implicit neural representations (INR) for continuous angular interpolation. We will consider integrating polarization light field imaging to ensure robust performance under extreme specularities, and strive to explore broader applications in dynamic 3D reconstruction and autonomous navigation.


\subsection{Discussion}

In this work, we have presented the Unified Dual-Mask Physical Model (UDMPM), fundamentally shifting light field (LF) depth estimation from heuristic, Lambertian-dependent patchworks to a physics-driven, unified framework. By explicitly decoding complex light-matter interactions at the pixel level, our approach overcomes the long-standing bottleneck of single Lambertian assumptions, demonstrating robust performance across diverse reflectance profiles, including pure diffuse, specular, and scattering surfaces.

The core of this paradigm shift lies in the dual-mask modulation mechanism. Traditional Epipolar Plane Image (EPI) based methods inherently fail when the view-consistency assumption is violated by non-Lambertian reflections, often treating valuable specular or scattering cues as noise. By introducing the medium mask ($M_{med}$) to quantify pixel-level surface roughness ($r = a/\lambda$) and the angular direction mask ($M_{ang}$) to capture the wavevector deflection field ($\Delta k$), UDMPM provides a rigorous physical boundary for light-matter interactions. Furthermore, our integration of Radiance Tensor Field (RTF) rank analysis theoretically grounds this design: while Lambertian scenes yield an RTF rank of 1, the superposition of multi-roughness and deflection vectors in non-Lambertian regions elevates the rank to $\ge 2$. This mathematical decoupling allows the network to adaptively modulate epipolar geometries without discarding high-frequency physical cues.

From a signal processing perspective, the MRI-inspired angular frequency analysis offers a highly efficient and physically interpretable solution for material parsing. Mapping the sparse $9 \times 9$ angular samples to a k-space equivalent and applying 2D-DFT elegantly circumvents the spatial-domain ambiguity of micro-textures, where conventional CNNs often overfit without physical guarantees. The distinct frequency signatures—low-frequency dominance for diffuse, high-frequency broadening for specular, and mid-frequency concentration for scattering—enable an $\mathcal{O}(HW)$ complexity material classification. This lightweight first tier of our three-step inverse wavevector parsing pipeline ensures real-time applicability, seamlessly feeding into the subsequent least-squares BRDF estimation and full wavevector resolution. However, we acknowledge certain limitations. The current 2D-DFT formulation assumes local spatial stationarity within the angular aperture; in scenarios with severe occlusions or extreme depth discontinuities where the angular signal is heavily aliased, the frequency mapping may degrade. Future iterations could incorporate non-uniform Fourier transforms or learned sparse representations to enhance robustness under severe occlusion.

Beyond depth estimation, the physical interpretability of UDMPM holds profound implications for downstream computational imaging tasks. In autonomous driving and robotic manipulation, understanding not just the geometry but also the material properties (e.g., distinguishing a reflective puddle from a diffuse shadow) is critical for safe navigation. The explicit extraction of $M_{med}$ and $M_{ang}$ provides these material priors natively, bridging the gap between geometric reconstruction and semantic scene understanding. Future work will explore extending this unified physical framework to dynamic light fields and integrating more complex, anisotropic Bidirectional Reflectance Distribution Function (BRDF) models to capture highly directional scattering phenomena. Ultimately, UDMPM establishes a foundational blueprint for physics-aware neural rendering and 4D scene reconstruction.

\subsection{Limitations and Future Work}

While the proposed Unified Dual-Mask Physical Model (UDMPM) establishes a robust theoretical framework for non-Lambertian light field (LF) depth estimation, several limitations remain, which naturally dictate our future research directions from both signal processing and optical engineering perspectives.

First, the current MRI-inspired angular frequency analysis relies on a regular $9 \times 9$ angular sampling grid to perform 2D Discrete Fourier Transform (DFT). In scenarios with highly sparse or irregular angular sampling (e.g., $4 \times 4$ grids or wide-baseline multi-view setups), the angular frequency resolution degrades significantly. This causes spectral aliasing between high-frequency specular lobes and mid-frequency scattering components, potentially compromising the pixel-wise material classification and the accuracy of the medium mask ($M_{med}$). Furthermore, the current 2D-DFT assumes spatial stationarity within the local angular patch, which may introduce boundary artifacts at sharp material transitions. Future work will explore Non-Uniform Fast Fourier Transform (NUFFT) and graph signal processing techniques in the angular domain to achieve robust spectral parsing under arbitrary sampling topologies.

Second, although the dual-mask mechanism successfully unifies isotropic diffuse, specular, and scattering interactions via scalar surface roughness and 2D wavevector shifts, it inherently simplifies the high-dimensional Radiance Tensor Field (RTF). Complex material properties, such as anisotropic reflectance (e.g., brushed metals, hair) and sub-surface scattering (e.g., human skin, translucent plastics), cannot be fully parameterized by the current isotropic scalar and vector masks. To address this, we plan to extend the dual masks into high-order tensor representations and incorporate polarized light field cues. This extension will enable the physical decoding of anisotropic Bidirectional Reflectance Distribution Functions (BRDFs) and volumetric scattering, further generalizing the unified physical model to extremely complex real-world surfaces.

Finally, from an engineering perspective, the full inverse wavevector parsing and dual-mask modulation incur substantial computational and memory overhead when scaled to ultra-high-resolution LF data or dense video streams. Future efforts will focus on hardware-software co-design to alleviate this bottleneck. By integrating emerging computational imaging sensors, such as polarization-sensitive metasurfaces or event-based plenoptic cameras, we aim to directly capture high-frequency angular and photometric cues at the optical frontend. This would fundamentally bypass the spatial-angular resolution trade-off inherent in conventional microlens arrays, facilitating real-time, physically grounded 3D perception for next-generation autonomous navigation and robotic manipulation systems.